\section{ماژول admm\_compressor (ADMM Joint Pruning and Quantization Compressor)}
\label{sec:module_admm_compressor}

\subsection{معرفی و انگیزه}
این مکانیزم بر مبنای مبانی نظری پیشرفته در ادبیات توجه و فشرده‌سازی \cite{ref_admmcompressor_2024} توسعه یافته است.

فشرده‌سازی مدل‌های یادگیری عمیق معمولاً در قالب دو فرآیند مجزا شامل هرس کردن وزن‌ها \lr{(Pruning)} و کوانتیزاسیون مقادیر \lr{(Quantization)} دنبال می‌شود. اجرای ترتیبی و گسسته این دو گام اغلب منجر به انباشت خطای تقریب و افت فاحش عملکرد مدل در نرخ‌های فشرده‌سازی بالا می‌گردد؛ چرا که هر گام بدون اطلاع از محدودیت‌های گام دیگر عمل می‌کند.

ماژول فشرده‌ساز \lr{ADMMCompressor} با به کارگیری روش جهت‌های متناوب ضرایب چندگانه لاگرانژ \lr{(Alternating Direction Method of Multipliers - ADMM)}، چارچوبی یکپارچه و بهینه‌سازی‌شده برای فشرده‌سازی توأم وزن‌ها ارائه می‌دهد. در این پارادایم، محدودیت‌های ساختاری تنک‌سازی (نظیر الگوهای سخت‌افزاری \lr{2:4 Structured Sparsity}) و محدودیت‌های کوانتیدگی یکنواخت (نظیر کوانتیزاسیون ۴بیتی به ازای هر کانال) به طور هم‌زمان در قالب متغیرهای کمکی فرمول‌بندی شده و تعادل بهینه‌ای میان حفظ دقت، کاهش فضای دیسک و شتاب‌بخشی استنتاج روی شتاب‌دهنده‌های مدرن تنسور کُر \lr{(Tensor Cores)} به دست می‌آید.

\subsection{مبانی نظری و فرمول‌بندی ریاضی}
مسئله بهینه‌سازی تلفیقی برای هر لایه خطی با وزن اولیه $\mathbf{W}_0$ به صورت زیر مدل‌سازی می‌گردد:
\begin{equation}
\min_{\mathbf{W}, \mathbf{Z}, \mathbf{Q}} \frac{1}{2} \|\mathbf{W} - \mathbf{W}_0\|_F^2 + g(\mathbf{Z}) + h(\mathbf{Q}) \quad \text{s.t.} \quad \mathbf{W} = \mathbf{Z}, \quad \mathbf{W} = \mathbf{Q}
\end{equation}
که در آن $\mathbf{Z}$ متغیر کمکی تنک، $\mathbf{Q}$ متغیر کمکی کوانتیده، $g(\mathbf{Z})$ تابع مشخصه مجموعه تنک (با نرخ تنک‌سازی هدف یا الگوی ۲:۴) و $h(\mathbf{Q})$ تابع مشخصه شبکه کوانتیزاسیون گسسته است. تابع لاگرانژی افزوده \lr{(Augmented Lagrangian)} متناظر عبارت است از:
\begin{equation}
\mathcal{L}_{\rho_1, \rho_2} = \frac{1}{2} \|\mathbf{W} - \mathbf{W}_0\|_F^2 + \frac{\rho_1}{2} \|\mathbf{W} - \mathbf{Z} + \mathbf{U}_1\|_F^2 + \frac{\rho_2}{2} \|\mathbf{W} - \mathbf{Q} + \mathbf{U}_2\|_F^2 + g(\mathbf{Z}) + h(\mathbf{Q})
\end{equation}
که در آن $\mathbf{U}_1$ و $\mathbf{U}_2$ متغیرهای دوگان مقیاس‌شده و $\rho_1, \rho_2 > 0$ ضرایب جریمه هستند. الگوریتم در هر گام تکرار $k$ مراحل متناوب زیر را طی می‌کند:

\paragraph{۱. به‌روزرسانی بسته وزن‌های پیوسته ($\mathbf{W}$):}
با صفر قرار دادن گرادیان نسبت به $\mathbf{W}$، حل بسته و یکتای زیر حاصل می‌شود:
\begin{equation}
\mathbf{W}^{(k+1)} = \frac{\mathbf{W}_0 + \rho_1 (\mathbf{Z}^{(k)} - \mathbf{U}_1^{(k)}) + \rho_2 (\mathbf{Q}^{(k)} - \mathbf{U}_2^{(k)})}{1 + \rho_1 + \rho_2}
\end{equation}

\paragraph{۲. تصویر بر مجموعه تنک ($\mathbf{Z}$):}
متغیر $\mathbf{Z}$ با تصویر بر مجموعه قید $g$ به‌روزرسانی می‌گردد؛ در حالت تنک‌سازی ساختاریافته ۲:۴، از هر ۴ درایه متوالی در سطرها، ۲ درایه با کمترین قدرمطلق صفر می‌شوند:
\begin{equation}
\mathbf{Z}^{(k+1)} = \Pi_{\mathcal{S}}\left(\mathbf{W}^{(k+1)} + \mathbf{U}_1^{(k)}\right)
\end{equation}

\paragraph{۳. تصویر بر شبکه کوانتیزاسیون یکنواخت ($\mathbf{Q}$):}
نگاشت به نزدیک‌ترین نقطه شبکه کوانتیده به صورت به ازای هر کانال انجام می‌پذیرد:
\begin{equation}
\mathbf{Q}^{(k+1)} = \Pi_{\mathcal{Q}}\left(\mathbf{W}^{(k+1)} + \mathbf{U}_2^{(k)}\right) = \operatorname{clamp}\left(\operatorname{round}\left(\frac{\mathbf{W}^{(k+1)} + \mathbf{U}_2^{(k)}}{\mathbf{s}}\right), q_{\min}, q_{\max}\right) \odot \mathbf{s}
\end{equation}

\paragraph{۴. به‌روزرسانی متغیرهای دوگان ($\mathbf{U}_1, \mathbf{U}_2$):}
\begin{align}
\mathbf{U}_1^{(k+1)} &= \mathbf{U}_1^{(k)} + \mathbf{W}^{(k+1)} - \mathbf{Z}^{(k+1)} \\
\mathbf{U}_2^{(k+1)} &= \mathbf{U}_2^{(k)} + \mathbf{W}^{(k+1)} - \mathbf{Q}^{(k+1)}
\end{align}

\paragraph{تحلیل طیفی و نامساوی‌های پایداری خطا}
برای بررسی میزان دقت تقریب ماتریس وزن یا نگاشت توجه در ماژول \lr{admm\_compressor}، تفاضل ماتریس فاکتورگیری‌شده از نگاشت متراکم اصلی بر اساس نامساوی هافدینگ-فون‌نیومان ارزیابی می‌گردد. فرض کنید $\mathbf{W} \in \mathbb{R}^{d_{\text{in}} \times d_{\text{out}}}$ ماتریس وزن هدف باشد که توسط تقریب کم‌مرتبه یا پراکنده با رتبه موثر $r$ جایگزین شده است. کران بالای خطای بازسازی بر مبنای نرم فروبنیوس به صورت زیر تصریح می‌شود:
\begin{equation}
\label{eq:spectral_bound_admm_compressor}
\| \mathbf{W} - \mathbf{W}_{r} \|_F = \sqrt{\sum_{i=r+1}^{\min(d_{\text{in}}, d_{\text{out}})} \sigma_i^2(\mathbf{W})} \le \epsilon \| \mathbf{W} \|_F
\end{equation}
که در آن $\sigma_i(\mathbf{W})$ بیانگر $i$-امین مقدار منفرد مرتب‌شده ماتریس به صورت نزولی است. نسبت حفظ انرژی طیفی $\eta_r = \frac{\sum_{i=1}^r \sigma_i^2}{\sum_{i=1}^{d} \sigma_i^2}$ تضمین می‌کند که با انتخاب آستانه $\tau \ge 0.90$، بیش از نود درصد از پراکندگی واریانس داده‌ها حفظ شده و افت سیگنال در انتشار گرادیان رو به عقب به حداقل ممکن کاهش می‌یابد.

\paragraph{تحلیل مجانبی و پیچیدگی اشتراک حافظه سخت‌افزاری}
از منظر پیچیدگی محاسباتی و انتقال داده‌ها در سلسومراتب حافظه (حافظه سراسری \lr{HBM} به حافظه محلی \lr{SRAM})، عملیات ضرب ماتریسی در این مکانیزم نیازمند تخصیص بافر موقت با ابعاد فاکتورگیری‌شده است. برای پردازش دنباله‌ای به طول $N$ با بعد مخفی $D$، بار محاسباتی هسته از فرم درجه دو به رابطه زیر تقلیل می‌یابد:
\begin{equation}
\label{eq:complexity_admm_compressor}
\mathcal{C}_{\text{admm\_compressor}} = \mathcal{O}\left( B \cdot N \cdot D \cdot r + B \cdot H \cdot N \cdot d_k \right) \ll \mathcal{O}\left( B \cdot H \cdot N^2 \cdot d_k \right)
\end{equation}
این بهبود مرتبه پیچیدگی، با کاهش چشمگیر حجم تبادل تانسورها در گذرگاه‌های ارتباطی کارت گرافیک، مانع از اشباع پهنای باند حافظه (پدیده \lr{Memory-Bound}) در فرآیند استنتاج طولانی‌مدت می‌گردد.

\subsection{معماری و پیاده‌سازی}
کلاس اصلی پیاده‌سازی‌شده در فایل \lr{Codes/framework/compression\_engine/admm\_compressor.py} به نام \lr{ADMMCompressor} از نوآوری‌های مهندسی زیر بهره‌مند است:
\begin{itemize}
\item \textbf{کوانتیزاسیون مقاوم به ازای هر کانال \lr{(Per-Channel Quantization)}:} جلوگیری از خطای مقیاس ناشی از وزن‌های فرین \lr{(Outliers)} با تخصیص ضرایب مقیاس مجزا به ازای هر سطر خروجی ماتریس وزن.
\item \textbf{ردیابی بهترین تکرار بر پایه وزن واقعی W:} محاسبه دقیق خطای اعتبارسنجی بر روی ماتریس وزن اصلی و جلوگیری از ناپایداری نوسانی متغیرهای کمکی در تکرارهای پایانی.
\item \textbf{حفاظت از لایه‌های حیاتی \lr{(Critical Layer Protection)}:} مستثنی کردن خودکار لایه‌های طبقه‌بند، سر مدل \lr{(LM Head)} و نرمال‌سازی‌های لایه جهت جلوگیری از فروپاشی کل مدل.
\end{itemize}

\subsection{نتایج تجربی}
ارزیابی آزمایشگاهی ماژول \lr{admm\_compressor} بر روی چهار مدل مبنا نتایج قابل توجه مندرج در جدول~\ref{tab:benchmark_admm_compressor} را نشان می‌دهد:

\begin{table}[H]
\centering
\small
\caption{کارنامه تجربی و ارزیابی بنچ‌مارک سخت‌افزاری ماژول \lr{admm\_compressor} روی مدل‌های پایه}
\label{tab:benchmark_admm_compressor}
\begin{tabular}{lccccc}
\toprule
\textbf{مدل پایه \lr{(Base Model)}} & \textbf{دقت \lr{(Accuracy)}} & \textbf{شاخص \lr{F1}} & \textbf{حافظه \lr{(Peak VRAM)}} & \textbf{توان \lr{(Throughput)}} & \textbf{وضعیت} \\
\midrule
\lr{DistilBERT} & \lr{88.53\%} & \lr{0.8922} & \lr{339.42 MB} & \lr{54,607.47 tok/s} & \lr{PASS} \\
\lr{RoBERTa} & \lr{89.11\%} & \lr{0.8966} & \lr{333.44 MB} & \lr{13,290.18 tok/s} & \lr{PASS} \\
\lr{GPT-2} & \lr{84.63\%} & \lr{0.8537} & \lr{409.87 MB} & \lr{15,433.04 tok/s} & \lr{PASS} \\
\lr{Qwen} & \lr{88.42\%} & \lr{0.8947} & \lr{1,851.31 MB} & \lr{1,099.69 tok/s} & \lr{PASS} \\
\bottomrule
\end{tabular}
\end{table}

\subsection{محدودیت‌ها و تحلیل شکست}
تحلیل چندوجهی کارایی این ماژول نکات برجسته‌ای را آشکار می‌سازد:
\begin{itemize}
\item \textbf{پایداری همه‌جانبه با نمره قبولی کامل (PASS):} تمامی چهار مدل تحت آزمون موفق به کسب نرخ بالای دقت (بین \lr{84.63\%} تا \lr{89.11\%}) با همگرایی پایدار لاگرانژی شدند.
\item \textbf{تعادل قابل توجه دقت و مصرف حافظه:} مدل \lr{RoBERTa} با دقت \lr{89.11\%} و F1 معادل \lr{0.8966} و مصرف حافظه تنها \lr{333.44 MB}، گواهی بر قدرت همگرایی هم‌زمان هرس و کوانتیزاسیون است.
\item \textbf{چالش زمان کالیبراسیون و بار محاسباتی بهینه‌سازی:} علی‌رغم نتایج استنتاجی عالی، اجرای تکرارهای متناوب حل بسته و تصاویر غیرخطی نیازمند صرف زمان پردازشی اولیه و حافظه محاسباتی برای نگهداری متغیرهای دوگان ($\mathbf{U}_1, \mathbf{U}_2$) در طول فشرده‌سازی است. در صورتی که تعداد تکرارها کمتر از ۴ تنظیم شود، واگرایی دوگان رخ داده و خطای بازسازی به شدت افزایش می‌یابد.
\end{itemize}
